\section{Five factors and extraction at one place}
At a contact prime $p\notin\Sigma$, $\delta$ is a unit and
$\nu=v_p(R)=v_p(F_2(t_*))$ is positive and odd, and $v_p(q)=-\nu$. This follows
from the unique contact factor with positive odd valuation and
does not require any condition on the valuation of the scale. All valuations
of $K$ are unramified and normalized by $v(p)=1$.
For a point on $E$, every finite value of $v(f)$ is even: if
$v(x)<0$ then $2v(y)=3v(x)$, while if $v(x)>0$ then
$2v(y)=v(x)$; the remaining case is valuation zero.
This also holds for conjugated reciprocal abscissas. The common-point lemma identifies $F_0=f(H-P)$ on every solution,
including the exceptional leading-coefficient cases. Its even valuation
is therefore available before the exclusion of torsion classes and
before any formal-parameter quotient is used.

Use
\begin{align*}
 T_1&=1+q^5C_2^2+q^2D_H^2+q^4B_1^4,\\
 T_2&=1+q^5G^2+q^2A_3^2+q^4B_0^4,\\
 T_3&=1+q^5I^4+q^3A_2^2+q^2A_4^2+q^4B_3^4,\\
 T_4&=1+q^5C_3^2+q^2C_1^4+q^4B_4^4,\\
 T_5&=1+q^5C_4^2+q^2A_1^2+q^4B_2^4,
 \qquad W=T_1T_2T_3T_4T_5.
\end{align*}

\begin{lemma}\label{lem:residues}
At a contact prime, every factor is nonzero and has valuation at most
zero. A factor of valuation zero is congruent to one in the residue field.
\end{lemma}
\begin{proof}
First suppose every displayed test in a factor is nonzero. In each
factor other than $T_3$, the three nonconstant valuations belong,
respectively, to the odd integers, to $2+4\Z$, and to $4+8\Z$.
For example $T_1$ has valuations
\[
 2v(C_2)-5\nu,\qquad 2v(D_H)-2\nu,\qquad4v(B_1)-4\nu.
\]
The second has exact factor two because $v(D_H)$ is even and $\nu$ is
odd. The third has exact factor four because $v(B_1)$ is even.
The term $q^2C_1^4$ in $T_4$ has valuation
$2(2v(C_1)-\nu)$, again twice an odd integer.

In $T_3$ the four nonconstant valuations are
\[
 4v(I)-5\nu,\quad 2v(A_2)-3\nu,\quad
 2v(A_4)-2\nu,\quad4v(B_3)-4\nu.
\]
The first three have residues $-\nu,\nu,2$ modulo four. These are
pairwise distinct and nonzero because $\nu$ is odd and the $A_i$
have even valuations. The fourth has residue four modulo eight;
it therefore differs from each of the first three and from zero.
Thus, in every factor, all finite term valuations, including zero
for the constant term, are distinct. Zero tests simply remove terms
from this argument. The minimum is attained exactly once. This proves
nonvanishing and the asserted bound. If the minimum is zero, every
nonconstant term has positive valuation, so the residue is one.
\end{proof}

\begin{proposition}\label{prop:extract}
At each contact prime, the retained pole system and the cleared trace
equations imply $\nu>0$ and give all sixteen required test orders at
one place above that prime.
\end{proposition}
\begin{proof}
Write $W=w_0+w_1\rho$. The literal cleared trace equations are
\[
 w_0^2-2w_0\tau-2w_1^2=0,\qquad (1+q)\mathfrak b=\gamma^2\tau^2.
\]
By Lemma~\ref{lem:residues}, $W\ne0$. The first equation therefore
gives $w_0\ne0$ and $\tau=1/\Tr(W^{-1})\ne0$; indeed $w_0=0$
would give $w_1=0$. Hence
$\mathfrak b=\gamma^2R/\Tr(W^{-1})^2$. The factor $\gamma$ is a unit at every consumed contact prime. If
$\ell=\max_{v\mid p}v(W)\le0$, the trace inequality gives
$v_p(\mathfrak b)\le\nu+2\ell$.
This valuation is odd; the retained pole obstruction forces it positive.
Consequently $\nu>0$ and $\ell>-\nu/2$.
At a place attaining $\ell$, the factor valuations are nonpositive
and add to $\ell$. Each factor valuation is therefore at least $\ell$,
and each of its term valuations is strictly greater than $-\nu/2$.

The resulting lower bounds, before integer rounding, are
\[
\begin{array}{c|c|c}
\text{test}&\text{strict lower bound}&\text{integer order}\\\hline
C_2,C_3,C_4,G&9\nu/4&3\\
I&9\nu/8&2\\
C_1&3\nu/8&1\\
D_H&3\nu/4&2\\
A_2&5\nu/4&1\\
A_1,A_3,A_4&3\nu/4&1\\
B_0,B_1,B_2,B_3,B_4&7\nu/8&1
\end{array}
\]
Here $\nu$ is a positive odd integer. The order two for $D_H$ uses
its even valuation. In fact each of the nine reciprocal-abscissa
tests has order at least two; order one suffices for its subsequent
formal-group use. Zero tests satisfy all the stated inequalities.
\end{proof}

\section{Integrality before the formal congruences}
At a contact place selected by Proposition~\ref{prop:extract}, the
points $H,U_i$ are formal in the primary completion, and $H-P,V_i$
are formal in the conjugate completion. The reciprocal-abscissa
tests have positive even valuation, so every nonidentity point has
formal depth at least one. In the parameter $z=-x/y$, write
\[
 \log_E(Q)=z(Q)\mathcal A(Q),\qquad
 \mathcal A(Q)\equiv1\pmod{p^{4v(z(Q))}}.
\]
The value of the unit is taken to be one at the identity. This is
the constant-term formal logarithm estimate proved in~\cite{Logs}.

The integral group identities
\[
                  [m]U_i=[2u_i]H,\qquad
                  [m-2]V_i=[2(u_i-n_i)](H-P)
\]
kill the torsion in $H=[m]R_0+T_0$. Their full logarithmic ratios
are $u_i/h$ and $(u_i-n_i)/(h-1)$, where
$(n_1,n_2,n_3,n_4)=(n,a,t,b)$. At least one of $h,h-1$ is a
$p$-adic unit, so at least one full ratio is integral. The formal
units preserve the valuation of its plain parameter ratio.
The positive-order test $C_i$ transfers integrality to the other
plain ratio, and then to the other full ratio.

Only after this transfer may the relative errors of order four
be used as absolute congruences. They give
\[
 L_i\equiv u_i/h\pmod{p^4},\qquad
 L_i'\equiv(u_i-n_i)/(h-1)\pmod{p^4}.
\]
Since $n_i=h(u_i/h)-(h-1)((u_i-n_i)/(h-1))$, the high cross tests
give $L_i\equiv n_i\pmod{p^3}$ for $i>1$. The first cross test
gives $L_1\equiv n\pmod p$, while $I$ gives $L_1\equiv\beta\pmod p$.
These arguments include identity numerator points, half-integral $h$,
and inert primes, where conjugation preserves the unique valuation.

